\section*{Capítulo \ref{ch:g10:earth-energy-budget} --- O balanço de energia da Terra}

\begin{solution}{exo:g10:earth-energy-budget:1}
A \omterm{def:g10:earth-energy-budget:insolation}{insolação} de um lugar é a potência da luz solar recebida por um metro
quadrado horizontal no topo da atmosfera acima dele (em média sobre um dia
ou um ano). O \omterm{def:g10:earth-energy-budget:albedo}{albedo} de uma superfície é a fração da luz solar que incide
sobre ela que ela reflete. A neve fresca (0.8 a 0.9) reflete muito mais que
o oceano aberto (0.06).
\end{solution}

\begin{solution}{exo:g10:earth-energy-budget:2}
Refletida: $100/340 \approx \qty{29}{\percent}$; absorvida pelo ar:
$80/340 \approx \qty{24}{\percent}$; absorvida pela superfície: $160/340
\approx \qty{47}{\percent}$.
\end{solution}

\begin{solution}{exo:g10:earth-energy-budget:3}
Asfalto: $800 \times (1 - 0.04) = \qty{768}{W}$. Neve: $800 \times (1 -
0.85) = \qty{120}{W}$, mais de seis vezes menos.
\end{solution}

\begin{solution}{exo:g10:earth-energy-budget:4}
Evaporação: \qty{82}{W/m^2}, contra \qty{21}{W/m^2} para o aquecimento do
ar. A água que evapora retira o seu \omterm{def:g10:earth-energy-budget:latent-heat}{calor latente} do mar; o vapor sobe com
o ar e, onde se condensa em gotículas de nuvem, muitas vezes quilômetros
mais acima, devolve essa energia ao ar em volta.
\end{solution}

\begin{solution}{exo:g10:earth-energy-budget:5}
Em dezembro, o hemisfério norte se inclina para longe do Sol: no polo, o
Sol fica abaixo do horizonte o dia inteiro, logo nenhuma luz solar chega.
No solstício de dezembro, isso vale em toda parte, em direção ao polo, além de $90^\circ - 23.44^\circ =
\qty{66.56}{\degree}$ N.
\end{solution}

\begin{solution}{exo:g10:earth-energy-budget:6}
$1361 \sin 90^\circ = \qty{1361}{W}$; $1361 \sin 30^\circ =
\qty{680.5}{W}$; $1361 \sin 10^\circ \approx \qty{236}{W}$. A metade de
$S_0$ exige $\sin h = 0.5$, isto é, $h = \qty{30}{\degree}$.
\end{solution}

\begin{solution}{exo:g10:earth-energy-budget:7}
Refletidos 100, absorvidos pelo ar 80, absorvidos pela superfície 160:
$100 + 80 + 160 = 340$, toda a luz solar está contabilizada. Na superfície,
$185 = 25 + 160$: o que não é refletido é absorvido.
\end{solution}

\begin{solution}{exo:g10:earth-energy-budget:8}
Disco: $\pi \times (6.371\times10^6)^2 \approx \qty{1.275e14}{m^2}$;
potência: $1361 \times 1.275\times10^{14} \approx \qty{1.74e17}{W}$.
Esfera: $4\pi
R^2 \approx \qty{5.10e14}{m^2}$. Média: $1.74\times10^{17}/5.10\times
10^{14} \approx \qty{340}{W/m^2}$, um quarto de $S_0$.
\end{solution}

\begin{solution}{exo:g10:earth-energy-budget:9}
$0.7 \times 5.10\times10^{14} \times 3.16\times10^7 \approx
\qty{1.1e22}{J}$ em um ano.
\end{solution}

\begin{solution}{exo:g10:earth-energy-budget:10}
As nuvens refletem $100 - 53 = \qty{47}{W/m^2}$ a mais de luz solar, que
portanto não é absorvida. Elas retêm $267 - 239 = \qty{28}{W/m^2}$ de
infravermelho. No saldo, a Terra perde $47 - 28 = \qty{19}{W/m^2}$: as
nuvens, como um todo, a resfriam.
\end{solution}

\begin{solution}{exo:g10:earth-energy-budget:11}
De cerca de 10 de maio até o fim de julho, isto é, em maio, junho e julho.
No equador, o Sol passa a pino ao meio-dia nos dois equinócios, em março e
em setembro; nos solstícios, ele passa a \qty{23.44}{\degree} do zênite,
logo o equador recebe um pouco menos em junho e em dezembro.
\end{solution}

\begin{solution}{exo:g10:earth-energy-budget:12}
A cada segundo: $82/(2.47\times10^6) \approx \qty{3.32e-5}{kg}$. A cada dia:
$3.32\times10^{-5} \times 86\,400 \approx \qty{2.9}{kg}$, uma camada de
cerca de \qty{2.9}{mm}. A cada ano: $3.32\times10^{-5} \times 3.16\times10^7 \approx
\qty{1050}{kg}$, cerca de \qty{1.05}{m} de água.
\end{solution}

\begin{solution}{exo:g10:earth-energy-budget:13}
$S_0 \sin\delta = S_0/\pi$ dá $\sin\delta = 1/\pi \approx 0.318$, logo
$\delta \approx \qty{18.6}{\degree}$. Então $\sin\lambda = 0.318/\sin 23.44^\circ
\approx 0.318/0.398 \approx 0.800$: $\lambda \approx \qty{53.15}{\degree}$ ou
$180^\circ - 53.15^\circ = \qty{126.85}{\degree}$. Entre os dois, $\lambda$
cresce \qty{73.7}{\degree}, isto é, $73.7/360$ de ano, cerca de
\num{75} dias, de meados de maio ao fim de julho.
\end{solution}

\begin{solution}{exo:g10:earth-energy-budget:14}
O gelo absorve $1 - 0.6 = 0.40$ da luz solar, o oceano $1 - 0.06 =
0.94$. Potência a mais: $(0.94 - 0.40) \times 300 \times 10^{12} \approx
\qty{1.6e14}{W}$, mais que toda a fotossíntese do planeta. A energia a mais
aquece a água, que derrete mais gelo nas bordas dos blocos, o que expõe mais
água escura: a mudança se alimenta a si mesma.
\end{solution}

\begin{solution}{exo:g10:earth-energy-budget:15}
Ponto mais próximo: $1361 \times (149.6/147.1)^2 \approx \qty{1408}{W/m^2}$;
mais afastado: $1361 \times (149.6/152.1)^2 \approx \qty{1317}{W/m^2}$;
cerca de \qty{7}{\percent} de diferença. A \qty{60}{\degree} N, a \omterm{def:g10:earth-energy-budget:insolation}{insolação}
é vinte vezes maior em junho do que em dezembro: uma variação de alguns por
cento não pode causar isso. Além disso, a Terra fica mais perto do Sol no
início de janeiro, durante o inverno do hemisfério norte, e os dois
hemisférios, à mesma distância do Sol, têm estações opostas: é a
inclinação do eixo, e não a distância, que faz as estações.
\end{solution}

\begin{solution}{pb:g10:earth-energy-budget:1}
\textbf{1.} $\pi \times (6.371\times10^6)^2 \approx \qty{1.275e14}{m^2}$.

\textbf{2.} $1361 \times 1.275\times10^{14} \approx \qty{1.735e17}{W}$.

\textbf{3.} $4\pi R^2 \approx \qty{5.101e14}{m^2}$.

\textbf{4.} $1.735\times10^{17}/5.101\times10^{14} \approx \qty{340}{W/m^2}$,
isto é, $S_0/4$.

\textbf{5.} Equador: \qty{1361}{W}; \qty{60}{\degree} N: $1361 \sin
30^\circ = \qty{680.5}{W}$.

\textbf{6.} $100/340 \approx 0.294$.

\textbf{7.} $340 - 100 = \qty{240}{W/m^2}$; sobre a Terra, $240 \times
5.101\times10^{14} \approx \qty{1.224e17}{W}$.

\textbf{8.} $53/340 \approx 0.156$. As nuvens aumentam o \omterm{def:g10:earth-energy-budget:albedo}{albedo} por um fator
$100/53 \approx 1.9$: elas quase o dobram.

\textbf{9.} $(240/5.67\times10^{-8})^{1/4} \approx \qty{255}{K}$, cerca de
\qty{33}{K} abaixo dos \qty{288}{K} da superfície.

\textbf{10.} $160 + 342 = \qty{502}{W/m^2}$.

\textbf{11.} $82 + 21 + 398 = \qty{501}{W/m^2}$. A superfície ganha cerca de
\qty{1}{W/m^2}: o desequilíbrio de \qty{0.7}{W/m^2}, arredondado.

\textbf{12.} $342/160 \approx 2.1$. O ar envia à superfície, como
infravermelho, mais que o dobro da energia da luz solar que ela absorve: é
o efeito estufa.

\textbf{13.} $82 \times 3.156\times10^7/(2.47\times10^6) \approx
\qty{1050}{kg}$, uma camada de cerca de \qty{1.05}{m} de espessura.

\textbf{14.} $0.7 \times 5.101\times10^{14} \approx \qty{3.6e14}{W}$.

\textbf{15.} $2805/72.07 \approx \qty{38.9}{kJ}$ por grama de carbono.

\textbf{16.} $104.9\times10^9$ toneladas são $1.049\times10^{17}$~g;
$1.049\times10^{17} \times 38.9\times10^3 \approx \qty{4.08e21}{J}$.

\textbf{17.} $4.08\times10^{21}/3.156\times10^7 \approx \qty{1.29e14}{W}$.

\textbf{18.} $1.29\times10^{14}/5.101\times10^{14} \approx
\qty{0.25}{W/m^2}$.

\textbf{19.} $0.7/0.25 \approx 2.8$: o desequilíbrio da Terra é quase três
vezes a potência que a fotossíntese armazena.

\textbf{20.} $1.29\times10^{14}/1.224\times10^{17} \approx 0.0011$: a
fotossíntese armazena \textbf{cerca de \qty{0.1}{\percent}} da luz solar que
a Terra absorve --- um milésimo.
\end{solution}
